\subsection{Preesquemas geométricamente irreducibles y geométricamente
conexos.}
\label{subsection:IV.2.4.5-fr}

\begin{proposition}[4.5.1]
\label{IV.2.4.5.1-fr}
Sean $k$ un cuerpo, $X$ un $k$-preesquema, $\Omega$ una extensión
algebraicamente cerrada de $k$. Sea $n$ \textup{(resp.} $n'$\textup{)} el cardinal del conjunto
de las componentes irreducibles \textup{(resp.} de las componentes conexas\textup{)} de
$X\otimes_k\Omega$; este número es independiente de la extensión
algebraicamente cerrada $\Omega$ elegida. Además, para toda extensión $K$ de
$k$, el cardinal del conjunto de las componentes irreducibles \textup{(resp.} de las
componentes conexas\textup{)} de $X\otimes_k K$ es $\leqslant n$ \textup{(resp.}
$\leqslant n'$\textup{)}.
\end{proposition}

\begin{proof}
En efecto, dos extensiones algebraicamente cerradas $\Omega$, $\Omega'$ de
$k$ siempre pueden considerarse como subextensiones de una misma extensión
$E$ de $k$; basta entonces aplicar (4.4.6) a $X\otimes_k\Omega$ y
$X\otimes_k\Omega'$ y a la extensión común $E$ de $\Omega$ y $\Omega'$ para
demostrar la primera afirmación. La segunda se obtiene tomando como $\Omega$
una extensión algebraicamente cerrada de $K$ y utilizando (4.4.1).
\end{proof}

\begin{definition}[4.5.2]
\label{IV.2.4.5.2-fr}
El cardinal $n$ \textup{(resp.} $n'$\textup{)} de \textup{(4.5.1)} se llama número geométrico de
componentes irreducibles \textup{(resp.} de componentes conexas\textup{)} de $X$
(relativamente a $k$). Si $n=1$ \textup{(resp.} $n'\leqslant1$\textup{)}, se dice que $X$ es un
$k$-preesquema geométricamente irreducible \textup{(resp.}
geométricamente conexo\textup{)}.
\end{definition}

Se recupera así una definición dada para un caso particular en (\textbf{III}, 4.3.4).
En virtud de (4.5.1), las propiedades siguientes son equivalentes:
\begin{enumerate}
\item[$a)$] $X$ es geométricamente irreducible (resp. geométricamente
conexo).
\item[$b)$] Para una extensión algebraicamente cerrada $\Omega$ de $k$,
$X\otimes_k\Omega$ es irreducible (resp. conexo).
\item[$c)$] Para toda extensión $K$ de $k$, $X\otimes_k K$ es irreducible
(resp. conexo).
\end{enumerate}

\begin{env}[4.5.3]
\label{IV.2.4.5.3-fr}
Sean $X$ un $k$-preesquema, $Z$ una parte localmente cerrada de $X$. Si se
denota también por $Z$ cualquiera de los subpreesquemas de $X$ que tienen $Z$
por espacio subyacente (\textbf I, 5.2.1), se deduce de (\textbf I, 5.1.8) que, para una
extensión $K$ de $k$, el número de componentes irreducibles (resp. conexas) de
$Z\otimes_k K$ es independiente del subpreesquema de espacio subyacente $Z$
que se haya elegido; por ello se define el número geométrico de componentes
irreducibles (resp. de componentes conexas) de $Z$ como el de cualquier
subpreesquema de $X$ que tenga $Z$ por espacio subyacente.
\end{env}

\begin{proposition}[4.5.4]
\label{IV.2.4.5.4-fr}
Sean $X$, $Y$ dos $k$-preesquemas, $f:X\to Y$ un $k$-morfismo sobreyectivo.
Si $X$ es geométricamente irreducible \textup{(resp.} geométricamente conexo\textup{)}, lo mismo
ocurre con $Y$.
\end{proposition}

\oldpage[IV-2]{62}

\begin{proof}
En efecto, para toda extensión $K$ de $k$, $f_{(K)}:X_{(K)}\to Y_{(K)}$ es
sobreyectivo (\textbf I, 3.5.2), y, siendo $X_{(K)}$ irreducible (resp. conexo) por
hipótesis, lo mismo ocurre con $Y_{(K)}$.
\end{proof}

\begin{definition}[4.5.5]
\label{IV.2.4.5.5-fr}
Se dice que un morfismo $f:X\to Y$ de preesquemas es irreducible
\textup{(resp.} conexo\textup{)} si, para todo $y\in Y$, el $\boldsymbol{k}(y)$-preesquema
$f^{-1}(y)$ es geométricamente irreducible \textup{(resp.} geométricamente conexo\textup{)}.
\end{definition}

\begin{proposition}[4.5.6]
\label{IV.2.4.5.6-fr}
\begin{enumerate}
\item[\textup{(i)}] Sean $X$ un $k$-preesquema, $K$ una extensión de $k$. Entonces el
número geométrico de componentes irreducibles \textup{(resp.} conexas\textup{)} de $X_{(K)}$
relativamente a $K$ es igual al número geométrico de componentes irreducibles
\textup{(resp.} conexas\textup{)} de $X$ relativamente a $k$. En particular, para que el
$K$-preesquema $X_{(K)}$ sea geométricamente irreducible \textup{(resp.}
geométricamente conexo\textup{)}, es necesario y suficiente que lo sea el
$k$-preesquema $X$.
\item[\textup{(ii)}] Sean $f:X\to Y$, $g:Y'\to Y$ dos morfismos, y pongamos
$X'=X_{(Y')}=X\times_Y Y'$, $f'=f_{(Y')}:X'\to Y'$. Si $f$ es irreducible
\textup{(resp.} conexo\textup{)}, lo mismo ocurre con $f'$. El recíproco es cierto si $g$ es
sobreyectivo.
\end{enumerate}
\end{proposition}

\begin{proof}
(i) Si $\Omega$ es una extensión algebraicamente cerrada de $K$, se
tiene $X\otimes_k\Omega=X_{(K)}\otimes_K\Omega$, de donde se sigue la
afirmación.

(ii) Para todo $y'\in Y'$, si se pone $y=g(y')$, se tiene
$f'^{-1}(y')=f^{-1}(y)\otimes_{\mathbf{k}(y)}\boldsymbol{k}(y')$ (\textbf I, 3.6.4), por lo que las dos
afirmaciones se deducen de (i).
\end{proof}

\begin{proposition}[4.5.7]
\label{IV.2.4.5.7-fr}
Sean $f:X\to Y$ un morfismo sobreyectivo irreducible \textup{(resp.} conexo\textup{)},
$g:Y'\to Y$ un morfismo, y pongamos $X'=X\times_Y Y'$. Si además $Y'$ es
irreducible \textup{(resp.} conexo\textup{)}, y $f$ es universalmente abierto \textup{(resp.} plano y
cuasicompacto, o universalmente abierto, o universalmente cerrado\textup{)}, entonces
$X'$ es irreducible \textup{(resp.} conexo\textup{)}.
\end{proposition}

\begin{proof}
Como todas las propiedades de $f$ consideradas son estables por cambio de
base, es posible limitarse al caso en que $Y'=Y$; basta entonces aplicar (4.4.2)
(teniendo en cuenta (2.3.12)).
\end{proof}

\begin{corollary}[4.5.8]
\label{IV.2.4.5.8-fr}
Sean $X$, $Y$ dos $k$-preesquemas.
\begin{enumerate}
\item[\textup{(i)}] Si $X$ es geométricamente irreducible \textup{(resp.} geométricamente
conexo\textup{)} e $Y$ irreducible \textup{(resp.} conexo\textup{)}, entonces $X\times_k Y$ es
irreducible \textup{(resp.} conexo\textup{)}.
\item[\textup{(ii)}] Si $X$ e $Y$ son geométricamente irreducibles \textup{(resp.}
geométricamente conexos\textup{)}, lo mismo ocurre con $X\times_k Y$.
\end{enumerate}
\end{corollary}

\begin{proof}
(i) El morfismo estructural $X\to\operatorname{Spec}(k)$ es sobreyectivo
y universalmente abierto (2.4.9), y además irreducible (resp. conexo); basta,
pues, aplicar (4.5.7).

(ii) Si $\Omega$ es una extensión algebraicamente cerrada de $k$, se
tiene $(X\times_k Y)_{(\Omega)}=X_{(\Omega)}\times_\Omega Y_{(\Omega)}$
(\textbf I, 3.3.10); por hipótesis, $X_{(\Omega)}$ e $Y_{(\Omega)}$ son
$\Omega$-preesquemas geométricamente irreducibles (resp. geométricamente
conexos) (4.5.6), y basta entonces aplicar (i).
\end{proof}
